\chapter{Estructura del repositorio}

La parte relevante del repositorio para esta memoria es:

\begin{verbatim}
phase1-pablo/
  src/decisionlab/
    agents/
      researcher.py
      formalizer.py
      reasoner.py
      builder.py
      memory_agent.py
    knowledge/
      extraction.py
      kg_writer.py
      indexer.py
      resolver.py
      retrieval/
        tool.py
        vector_retrieval.py
        kg_retrieval.py
        fusion.py
        crag.py
    runtime/
      loop.py
      dispatcher.py
    router.py
  docs/
    formal-documentation/
    tfg-memoria-latex/

shared/
  shared/
    services.py
    models.py
    storage.py
    database.py
    knowledge_graph.py
    vector_store.py
    pipeline_memories.py
\end{verbatim}

\section{Documentación formal}

La carpeta \texttt{phase1-pablo/docs/formal-documentation} contiene la
documentación técnica que sirve como base para esta memoria: arquitectura,
workflow de agentes, memoria, infraestructura compartida, decisiones de diseño y
bibliografía técnica.

\section{Comandos de referencia}

Para ejecutar pruebas de Python se usa:

\begin{verbatim}
uv run pytest
\end{verbatim}

Para compilar la memoria, desde la carpeta \texttt{tfg-memoria-latex}:

\begin{verbatim}
pdflatex traballo.tex
pdflatex traballo.tex
\end{verbatim}
